\chapter*{Resumen ejecutivo}
\addcontentsline{toc}{chapter}{Resumen ejecutivo}

\textbf{De cada 100 pesos que el Gobierno se propone recibir en 2026, 14,45 provienen de
deuda.} Eran 13,40 en la iniciativa anterior. El gasto neto total propuesto es de
10\,193\,683,7 millones de pesos, 26,3\,\% del PIB, y crece 4,6\,\% real contra el proyecto de
2025.

\textbf{Marco macroeconómico.} Crecimiento supuesto de 1,8 a 2,8\,\% después de cerrar 2025
entre 0,5 y 1,5. El PIB nominal estimado de 2025 quedó \emph{por debajo} del aprobado. El
supuesto más cargado no es el crecimiento sino la tasa: Cetes nominal promedio cae de 8,4 a
6,6\,\%.

\textbf{Ingresos.} La recaudación tributaria alcanza 15,1\,\% del PIB, el nivel más alto de la
serie reciente, \emph{sin reforma de tasas}: la miscelánea son tres leyes ---Derechos, IEPS y
Código Fiscal--- y no incluye iniciativa de ISR ni de IVA. La mitad del aumento real viene de
dos renglones que crecen sesenta por ciento: comercio exterior y accesorios de impuestos.

\textbf{Energía.} La transferencia del Fondo Mexicano del Petróleo cae 20,7\,\% real, sobre
todo por volumen exportado. La contrapartida está en el gasto: el Ramo 18 casi se duplica y
263\,476,3 de sus 267\,439,1 millones son una sola partida de compra de títulos y valores. El
Estado recibe menos renta y aporta más capital.

\textbf{Agregados.} 271 claves de programa aparecen y 358 desaparecen. Cinco ramos se
extinguen y tres nacen. El gasto corriente estructural queda 356\,980,1 millones por debajo de
su límite máximo: la regla fiscal que el paquete sí cumple.

\textbf{Salud.} La función crece 4,6\,\% real después de caer 12,2 el ejercicio anterior.
IMSS-Bienestar estrena ramo propio, el tercero en tres años. El fondo federalizado de salud no
se recuperó: se estabilizó en su nivel bajo. Por habitante, 7\,249 pesos al año.

\textbf{Educación.} Crece 3,4\,\% real; 2,6 por habitante y 4,4 por persona en edad escolar,
porque hay menos niños. El Ramo 11 crece 8,5\,\% y el federalizado 0,1. La media superior es el
único nivel que pierde.

\textbf{Pensiones.} \textbf{El agregado de pensiones y jubilaciones de la clasificación
económica cae 0,7\,\% real, por primera vez en la serie}, mientras las pensiones no
contributivas crecen 12,1\,\%. El mismo agregado cae 4,8\,\% por persona de 65 años y más. El
bloque contributivo absorbe el ajuste.

\textbf{Inversión.} Física, 2,5\,\% del PIB; financiera, 0,7, el punto más alto del horizonte,
que vuelve a 0,3 en 2027. Dentro de los capítulos 6000 y 7000 del objeto del gasto,
cerca de la mitad corresponde a inversiones financieras y provisiones, con la aportación de
capital a Pemex como componente dominante.

\textbf{Seguridad.} El perímetro de quince subfunciones queda plano, $-0,6$\,\% real. Debajo,
la Guardia Nacional pasó del Ramo 36 al Ramo 07. El gasto no cambia; cambia quién manda.

\textbf{Gasto federalizado.} Participaciones y aportaciones suman 6,45\,\% del PIB. Los
recursos de libre disposición crecen más rápido que los condicionados: 3,7\,\% real las
participaciones contra 1,5 las aportaciones, que también aumentan.

\textbf{Medio ambiente y agua.} 34\,535,3 millones, nueve centésimas de punto del PIB, con una
caída de 10,2\,\% real. El saneamiento de aguas residuales se reduce a la mitad.

\textbf{Anexos transversales.} Los doce suman el 74,3\,\% del gasto programable y no son
dinero adicional: 170 de 288 programas están en más de un anexo. El 48,1\,\% del anexo de
igualdad entre mujeres y hombres son programas de pensión.

\textbf{Deuda.} El saldo histórico sube de 51,4 a 52,3\,\% del PIB y el escenario oficial lo
deja ahí seis años. \emph{El paquete anterior prometió estabilidad en 51,4 y el acervo subió
nueve décimas en doce meses.} El costo financiero crece 8,1\,\% real y es la partida que más
crece.

\textbf{Horizonte demográfico.} El paquete proyecta pensiones a 2031 y no declara un solo
supuesto de población para educación ni para salud. Discute el horizonte donde la demografía
empuja al alza y calla donde empuja a la baja.
\clearpage
